%%%%%%%%%%%%%%%%%%%%%%%%%%%%%%%%%%%%%%%%%%%%%%%%%%%%%%%%%%%%%%%%%%%%%%%%
% Plantilla TFG/TFM
% Escuela Politécnica Superior de la Universidad de Alicante
% Realizado por: Jose Manuel Requena Plens
% Contacto: info@jmrplens.com / Telegram:@jmrplens
%%%%%%%%%%%%%%%%%%%%%%%%%%%%%%%%%%%%%%%%%%%%%%%%%%%%%%%%%%%%%%%%%%%%%%%%

\chapter{Desarrollo}
\label{desarrollo}

\section{Presentación de resultados}

Podría hacer un set de gráficas que muestren claramente la frontera de decisión para un lenguaje como 
anbn. Si un eje es el numero de a y el otro eje es el número de b, los modelos clasificatorios
deberían mostrar los valores predichos para cada par. Un modelo ideal debería mostrar 1 por toda la diagonal
y 0 por lo demás. Y mostrar lo mal que lo hace para valores de N hasta los cuales no ha podido generalizar 
todavía.
